\subsection{Clase GaloisField en GF(2)}
Para comenzar se instanció la clase previamente desarrollada para modelar campos de Galois con parametros $m = 1$ y $P(x) = x + 1$, para testear y corregir posibles issues. La implementación existente no requirió modificaciones y se validaron todas sus operaciones exitosamente:

\begin{figure}[H]
    \centering
    \includegraphics[scale=0.3]{gf2-test.png}
    \caption{Verificacion de las operaciones en GF2.}
    \label{fig:gf2-test}
\end{figure}

Una vez verificado el funcionamiento de la clase, se implementaron los metodos para realizar el producto vector-matriz,
matriz-matriz, peso de Hamming y algunos metodos auxiliares que se usaron en los modelos del encoder y decoder.

Usando estos nuevos metodos se verificaron las principales propiedades para las matrices del codigo ($GH^T = 0$ y $B^2 = I$):

\begin{figure}[H]
    \centering
    \includegraphics[scale=0.4]{code-matrices.png}
    \hfill
    \includegraphics[scale=0.4]{code-properties.png}
    \caption{Matrices y propiedades del código verificadas.}
    \label{fig:code-matrices-properties}
\end{figure}


Se genero ademas la tabla de síndromes que se encuentra en el archivo \href{https://github.com/AbelCorvalan0/fec-lab2/tree/main/model/outputs/syndromes/error_syndromes.txt}{\textbf{error syndromes.txt}}. Los vectores se presentan de la siguiente forma:

\begin{table}[h!]
    \centering
    \begin{tabular}{|c|c|}
        \hline
        1. Vector recibido     &
        2. Sindrome \\
        \hline
    \end{tabular}
    % \caption{Descripción de las columnas de salida del decodificador.}
    \label{tab:decoder_output}
\end{table}

% La totalidad de los errores que se pueden producir dentro de los 24 bits
% de una palabra recibida en el decoder son generados de la siguiente forma:

% El orden en que se encuentran las variables es el siguiente:

% \begin{table}[h!]
%     \centering
%     \begin{tabular}{|c|c|c|c|c|}
%         \hline
%         1. Vector recibido     &
%         2. Mensaje              &
%         3. Máscara de error     &
%         4. Palabra corregida    &
%         5. Palabra incorregible \\
%         \hline
%     \end{tabular}
%     \caption{Descripción de las columnas de salida del decodificador.}
%     \label{tab:decoder_output}
% \end{table}

% Los valores decodificados se encuentran en el archivo \href{https://github.com/AbelCorvalan0/fec-lab2/blob/main/model/outputs/encoder/encoder_top_testing/encoder_output.txt}{codeword with errors.txt}.
\subsection{Caracterización del decodificador}

Para cada peso entre 0 y 4, se recorre la totalidad de los patrones
de error correspondientes a dicho peso. Estos vectores son generados
en el golden model \href{https://github.com/AbelCorvalan0/fec-lab2/blob/main/model/golay_decoder.ipynb}{\textbf{golay decoder.ipynb}}.

Luego, estos vectores son procesados por un objeto decodificador para
obtener el archivo de vectores \href{https://github.com/AbelCorvalan0/fec-lab2/blob/main/model/outputs/decoder/decoder_top_testing/codeword_with_errors.txt}{\textbf{codeword with errors.txt}}.

Este archivo tiene el siguiente formato:
\begin{table}[h!]
    \centering
    \begin{tabular}{|c|c|c|c|c|}
        \hline
        1. Vector recibido &
        2. Mensaje &
        3. Máscara de error &
        4. Palabra corregida &
        5. Palabra incorregible \\
        \hline
    \end{tabular}
    %\caption{Descripción de las columnas de salida del decodificador.}
    \label{tab:decoder_output}
\end{table}

Con estos vectores se valido el modelo del decoder obteniendo exitosamente la tabla esperada de casos segun el peso de error:

\begin{figure}[H]
    \centering
    \includegraphics[scale=0.4]{decoder-error-table.png}
    \hfill
    \includegraphics[scale=0.4]{decoder-error-model.png}
    \caption{Tabla de casos de decodificacion segun peso de error (modelo).}
    \label{fig:decoder-error-table}
\end{figure}

Además se genera un conjunto de vectores por cada submódulo, todos los archivos se encuentran en \href{https://github.com/AbelCorvalan0/fec-lab2/tree/main/model/outputs/decoder}{\textbf{(vectores)}}.
